% Source: https://artofproblemsolving.com/community/c3h1374584_solve_the_equation
% Stable identity: aops-topic-1374584-q01
The equation requires $x\ne a$ and $x\ne-a$. Multiply both sides by $x^2-a^2=(x-a)(x+a)$, using $\dfrac{1}{a-x}=-\dfrac{1}{x-a}$:
\[
(3a+1)(x-a)+(a-1)(x+a)=2a(a^2-1).
\]
The left side is $3ax-3a^2+x-a+ax+a^2-x-a=4ax-2a^2-2a$, so
\[
4ax=2a^3-2a+2a^2+2a=2a^2(a+1).
\]
\begin{itemize}
\item If $a=0$: the equation reads $0=0$, so every $x$ with $x\ne\pm a$, i.e. every $x\ne0$, is a solution. (Directly: for $a=0$ the original equation is $\dfrac{1}{x}-\dfrac{-1}{-x}=0$, true for all $x\ne0$.)
\item If $a\ne0$: $x=\dfrac{a(a+1)}{2}$. This must differ from $\pm a$:
 $\dfrac{a(a+1)}{2}=a\iff a=1$, and $\dfrac{a(a+1)}{2}=-a\iff a=-3$. For these two values the only candidate is excluded, so there is no solution.
\end{itemize}
Summary:
\[
\boxed{x=\dfrac{a(a+1)}{2}\ \text{ if } a\ne0,1,-3;\quad\text{no solution if } a=1\text{ or }a=-3;\quad\text{every }x\ne0\text{ if }a=0.}
\]
Check with $a=2$: $x=3$, and $\dfrac{7}{5}-\dfrac{1}{-1}=\dfrac{12}{5}=\dfrac{2\cdot2\cdot3}{9-4}$.
